\documentclass[12pt, a4paper]{article}

% ── Encodage & langue ─────────────────────────────────────────────────────────
\usepackage[utf8]{inputenc}
\usepackage[T1]{fontenc}
\usepackage[french]{babel}

% ── Mathématiques ─────────────────────────────────────────────────────────────
\usepackage{amsmath, amssymb, amsthm}
\usepackage{mathtools}
\usepackage{bm}

% ── Mise en page ──────────────────────────────────────────────────────────────
\usepackage[margin=2.5cm]{geometry}
\usepackage{setspace}
\usepackage{parskip}
\usepackage{microtype}

% ── Figures & tableaux ────────────────────────────────────────────────────────
\usepackage{graphicx}
\usepackage{booktabs}
\usepackage{array}
\usepackage{float}

% ── Algorithmes ───────────────────────────────────────────────────────────────
\usepackage[ruled, vlined, linesnumbered]{algorithm2e}

% ── Hyperliens ────────────────────────────────────────────────────────────────
\usepackage[colorlinks=true, linkcolor=blue, citecolor=blue]{hyperref}

% ── Environnements mathématiques ──────────────────────────────────────────────
\theoremstyle{plain}
\newtheorem{theorem}{Théorème}[section]
\newtheorem{lemma}[theorem]{Lemme}
\newtheorem{proposition}[theorem]{Proposition}
\newtheorem{corollary}[theorem]{Corollaire}

\theoremstyle{definition}
\newtheorem{definition}[theorem]{Définition}
\newtheorem{remark}[theorem]{Remarque}

% ── Macros ────────────────────────────────────────────────────────────────────
\newcommand{\R}{\mathbb{R}}
\newcommand{\N}{\mathbb{N}}
\newcommand{\norm}[1]{\left\lVert #1 \right\rVert}
\newcommand{\inner}[2]{\left\langle #1,\, #2 \right\rangle}
\newcommand{\mat}[1]{\mathbf{#1}}
\newcommand{\vect}[1]{\mathbf{#1}}
\newcommand{\ones}{\mathbf{1}}
\newcommand{\Gram}{\mat{G}}
\newcommand{\Dsq}{\mat{D}^{\circ 2}}
\newcommand{\Jmat}{\mat{J}}
\newcommand{\tr}{\operatorname{tr}}

% ── En-tête ───────────────────────────────────────────────────────────────────
\title{
    \Large\textbf{Localisation de Microphones par\\
    Scaling Multidimensionnel et Trilatération}\\[0.5em]
    \large Application à l'Enregistrement Orchestral Multicanal
}
\author{
    Mehdi Maqlach\\[0.3em]
    \small Projet de spatialisation audio binaurale
}
\date{\today}

% ══════════════════════════════════════════════════════════════════════════════
\begin{document}
% ══════════════════════════════════════════════════════════════════════════════

\maketitle
\tableofcontents
\newpage

% ══════════════════════════════════════════════════════════════════════════════
\section{Introduction et formulation du problème}
% ══════════════════════════════════════════════════════════════════════════════

\subsection{Contexte}

Dans le cadre d'un système de spatialisation binaurale d'enregistrements
orchestraux, il est nécessaire de connaître la position tridimensionnelle de
chaque source sonore pour calculer les filtres HRTF appropriés. Le jeu de
données utilisé (TheOrchestra / TheSpheres, $f_s = 48\,000$ Hz) contient
$M = 23$ microphones disposés dans une salle de concert.

\subsection{Hypothèse de co-localisation}

Pour chaque instrument $i \in \{1, \ldots, 17\}$, il existe un microphone dit
\emph{dédié} $m^*_i$ placé en close-miking. On l'identifie par un critère
combinant \textbf{précocité} (faible $s_{ij}$) et \textbf{intensité} (fort
$a_{ij}$) :

\begin{equation}
    m^*_i \;=\; \underset{j \notin \{05\}}{\arg\min}\;\; \frac{s_{ij}}{a_{ij}}
    \label{eq:dedicated_mic}
\end{equation}

où $s_{ij} \in \N$ est le numéro d'échantillon du maximum absolu de la RIR
entre l'instrument $i$ et le microphone $j$, et $a_{ij} \in \R_{>0}$ est
l'amplitude (valeur absolue) de ce maximum.

\begin{remark}[Pourquoi le ratio $s_{ij}/a_{ij}$ et non $\arg\min\,s_{ij}$ seul]
    L'inspection des données révèle qu'un microphone éloigné peut présenter un
    échantillon $s_{ij}$ très faible (typiquement $s_{ij} \in \{1, 8, 10\}$)
    avec une amplitude $a_{ij}$ infime ($\sim 10^{-3}$). Ces détections
    précoces sont des artefacts de déclenchement ou des pré-échos du système
    d'acquisition, sans rapport avec le trajet direct acoustique. Le critère
    $\arg\min(s_{ij})$ seul serait alors trompé. Deux exemples illustratifs :

    \begin{center}
    \begin{tabular}{llccc}
        \toprule
        Instrument & Microphone & $s_{ij}$ & $a_{ij}$ & $s_{ij}/a_{ij}$ \\
        \midrule
        Basson & mic\textunderscore 18 (artefact) & 10 & 0{,}006 & 1724 \\
        Basson & mic\textunderscore 15 (dédié réel) & 33 & 0{,}404 & \textbf{82} \\
        \midrule
        Clarinette & mic\textunderscore 17 (artefact) & 10 & 0{,}004 & 2703 \\
        Clarinette & mic\textunderscore 14 (dédié réel) & 15 & 0{,}449 & \textbf{33} \\
        \bottomrule
    \end{tabular}
    \end{center}

    Le critère $\arg\min(s_{ij}/a_{ij})$ pénalise les artefacts précoces mais
    faibles, et sélectionne le microphone qui est à la fois \emph{proche}
    (latence faible) \emph{et} \emph{fort} (niveau élevé) — signature physique
    sans ambiguïté d'un microphone de proximité.
\end{remark}

\textbf{Hypothèse fondamentale.} On suppose que l'instrument $i$ est
positionné à la même position spatiale que son microphone dédié $m^*_i$.
Cette hypothèse est justifiée par le close-miking : le vrai dédié présente
systématiquement $s_{i,m^*_i} \approx 8$--$10$ échantillons \emph{et}
$a_{i,m^*_i} \gtrsim 0{,}2$, soit une distance source--microphone de l'ordre de
$5$--$7$ cm, négligeable devant les distances inter-microphones.

\subsection{Modèle de distance par TOA}

\subsubsection{Problème : maximum absolu \texorpdfstring{$\neq$}{≠} premier arrivé}

Le jeu de données fournit, pour chaque paire (instrument $i$, microphone $j$),
le numéro d'échantillon $s_{ij}$ du \textbf{maximum absolu} de la RIR, et non
celui de la \emph{première arrivée} (trajet direct). Dans une salle de concert
réverbérante, la RIR présente :

\begin{enumerate}
    \item le trajet direct, arrivant en premier mais souvent atténué par la distance,
    \item des réflexions précoces, potentiellement plus énergétiques que le direct,
    \item une queue de réverbération décroissante.
\end{enumerate}

Pour les paires de micros \emph{proches} (micro dédié co-localisé), le trajet
direct domine : $s_{i,m^*_i}$ est bien le TOA direct. Pour les paires
\emph{distantes}, une réflexion peut être le maximum de la RIR, donnant un
$s_{ij}$ bien supérieur au TOA direct et donc une distance surestimée.

\begin{remark}[Asymétrie révélatrice]
    La symétrie physique impose $d(m^*_i, m^*_j) = d(m^*_j, m^*_i)$, soit
    $s_{ij} \approx s_{ji}$.  Un écart important entre les deux mesures indique
    qu'une des deux a capturé une réflexion tardive plutôt que le trajet direct.
    Exemple dans les données :
    \begin{center}
    \begin{tabular}{lcc}
        \toprule
        Direction & $s$ (samples) & $d$ (m) \\
        \midrule
        Vln\textunderscore 1 $\to$ mic\textunderscore 18 (Trombones) & 2\,421 & 17{,}3 \\
        Trombones $\to$ mic\textunderscore 06 (Vln\textunderscore 1) &   838  &  6{,}0 \\
        \bottomrule
    \end{tabular}
    \end{center}
    L'écart d'un facteur 3 pour la même distance physique révèle que
    $s = 2\,421$ correspond à une réflexion murale tardive.
\end{remark}

\subsubsection{Correction : estimateur bidirectionnel}

Le trajet direct étant le plus court, il est borné supérieurement par le
minimum des deux mesures symétriques disponibles :

\begin{equation}
    d^{\rm direct}(m^*_i, m^*_j) \;\leq\;
    \frac{\min(s_{ij},\, s_{ji})}{f_s} \cdot c
    \;\leq\;
    \frac{\max(s_{ij},\, s_{ji})}{f_s} \cdot c
    \label{eq:toa_bounds}
\end{equation}

On adopte donc l'estimateur :

\begin{equation}
    \boxed{d_{ij} \;=\; \frac{\min(s_{ij},\, s_{ji})}{f_s} \cdot c}
    \label{eq:toa_distance}
\end{equation}

où $c = 343\,\text{m/s}$ et $f_s = 48\,000\,\text{Hz}$.

\begin{remark}[Propriété de symétrie]
    L'estimateur~\eqref{eq:toa_distance} est automatiquement symétrique :
    $d_{ij} = d_{ji}$.  La matrice $\mat{D}$ est donc symétrique par
    construction, sans nécessiter de post-traitement $(D + D^\top)/2$.
    Une asymétrie résiduelle nulle dans la matrice finale confirme que les
    deux sens de mesure convergent vers le même trajet le plus court.
\end{remark}

\subsection{Structure de la matrice de distances disponible}

Sous l'hypothèse de co-localisation, la mesure $d_{ij}$ est équivalente à la
distance entre le microphone dédié $m^*_i$ et le microphone $j$. On dispose
donc d'une matrice de distances $\mat{D} \in \R^{17 \times 23}$ dont la
structure induite sur la matrice complète $23 \times 23$ est la suivante.

On partitionne l'ensemble des 23 microphones en :
\begin{itemize}
    \item $\mathcal{D} = \{m^*_1, \ldots, m^*_{17}\}$ : les 17 microphones
        dédiés (mic\textunderscore 06 à mic\textunderscore 22),
    \item $\mathcal{A} = \{\text{mic\textunderscore 00}, \ldots,
        \text{mic\textunderscore 04}\}$ : les 5 microphones d'ambiance dont
        la distance à tout autre microphone d'ambiance est inconnue.
\end{itemize}

\begin{remark}
    Le microphone mic\textunderscore 05 présente des valeurs $s_{i,05}$
    anormalement élevées (entre 6\,387 et 30\,145 échantillons, soit des
    distances apparentes de 45 à 215 m physiquement impossibles dans une salle
    de concert). Ce microphone est exclu des calculs : il s'agit probablement
    d'une voie de réverbération artificielle ou d'un écho tardif.
\end{remark}

La matrice de distances complète $23 \times 23$ possède donc la structure
suivante, après symétrie par hypothèse de co-localisation :

\begin{equation}
    \mat{\Delta} \;=\; \begin{pmatrix}
        \mat{\Delta}_{\mathcal{DD}} & \mat{\Delta}_{\mathcal{DA}} \\[4pt]
        \mat{\Delta}_{\mathcal{AD}} & \mat{\Delta}_{\mathcal{AA}}
    \end{pmatrix}
\end{equation}

où $\mat{\Delta}_{\mathcal{DD}} \in \R^{17\times 17}$ et
$\mat{\Delta}_{\mathcal{DA}} \in \R^{17\times 5}$ sont \emph{connus}, mais
$\mat{\Delta}_{\mathcal{AA}} \in \R^{5\times 5}$ est \emph{inconnu} (aucun
instrument n'est co-localisé avec un microphone d'ambiance).

\medskip
Cette structure dicte une résolution en deux phases indépendantes :

\begin{enumerate}
    \item \textbf{Phase 1.} Retrouver les positions $\{\vect{x}_i\}_{i \in
        \mathcal{D}}$ des 17 microphones dédiés à partir de
        $\mat{\Delta}_{\mathcal{DD}}$ par \emph{Classical Multidimensional
        Scaling} (cMDS).
    \item \textbf{Phase 2.} Retrouver les positions $\{\vect{x}_k\}_{k \in
        \mathcal{A}}$ des 5 microphones d'ambiance à partir de
        $\mat{\Delta}_{\mathcal{AD}}$ par \emph{trilatération linéarisée}.
\end{enumerate}

% ══════════════════════════════════════════════════════════════════════════════
\section{Phase 1 : Classical MDS sur les microphones dédiés}
% ══════════════════════════════════════════════════════════════════════════════

\subsection{Objectif}

On dispose de la matrice de distances $\mat{\Delta}_{\mathcal{DD}} \in
\R^{17\times 17}$ symétrique à diagonale nulle. On cherche une configuration
$\mat{X} \in \R^{17 \times 3}$ dont les lignes sont les coordonnées
cartésiennes $\vect{x}_1, \ldots, \vect{x}_{17} \in \R^3$ telles que :

\begin{equation}
    \norm{\vect{x}_i - \vect{x}_j} \;=\; (\mat{\Delta}_{\mathcal{DD}})_{ij}
    \qquad \forall\, i, j \in \{1,\ldots,17\}
\end{equation}

\subsection{Pourquoi élever les distances au carré}

La distance euclidienne $d_{ij} = \norm{\vect{x}_i - \vect{x}_j}$ n'est pas
linéaire en les coordonnées. Son carré, en revanche, se décompose algébriquement :

\begin{align}
    d_{ij}^2 &= \norm{\vect{x}_i - \vect{x}_j}^2
    = (\vect{x}_i - \vect{x}_j)^\top (\vect{x}_i - \vect{x}_j) \notag \\
    &= \vect{x}_i^\top\vect{x}_i
       - 2\,\vect{x}_i^\top\vect{x}_j
       + \vect{x}_j^\top\vect{x}_j \notag \\
    &= \norm{\vect{x}_i}^2
       - 2\inner{\vect{x}_i}{\vect{x}_j}
       + \norm{\vect{x}_j}^2
    \label{eq:dist_squared}
\end{align}

Ce développement fait apparaître le \textbf{produit scalaire}
$\inner{\vect{x}_i}{\vect{x}_j}$, qui est une quantité \emph{bilinéaire} en
les coordonnées. C'est cette propriété que le MDS classique exploite : on
manipule une expression linéaire en les inconnues, ce qui est impossible avec
$d_{ij}$ directement.

\subsection{La matrice de Gram}

\begin{definition}[Matrice de Gram]
    Étant donné un ensemble de vecteurs $\vect{x}_1, \ldots, \vect{x}_n \in
    \R^d$, la matrice de Gram associée est :
    \begin{equation}
        \Gram \;=\; \mat{X}\mat{X}^\top \;\in\; \R^{n \times n},
        \qquad \Gram_{ij} = \inner{\vect{x}_i}{\vect{x}_j}
    \end{equation}
    où $\mat{X} \in \R^{n \times d}$ est la matrice dont les lignes sont les
    $\vect{x}_i$.
\end{definition}

L'objectif du MDS est de récupérer $\mat{X}$ (à une rotation près) à partir
de $\mat{G}$. L'équation~\eqref{eq:dist_squared} s'écrit alors sous forme
matricielle. En posant $r_i = \norm{\vect{x}_i}^2$ :

\begin{equation}
    \Dsq_{ij} \;=\; r_i + r_j - 2\,\Gram_{ij}
    \label{eq:gram_dist}
\end{equation}

ou, en notation matricielle :

\begin{equation}
    \Dsq \;=\; \vect{r}\ones^\top + \ones\vect{r}^\top - 2\Gram
    \label{eq:gram_dist_matrix}
\end{equation}

où $\vect{r} = (r_1, \ldots, r_n)^\top \in \R^n$ et $\ones = (1,\ldots,1)^\top
\in \R^n$.

\subsection{La matrice de centrage \texorpdfstring{$\Jmat$}{J}}

\begin{definition}[Matrice de centrage]
    Pour $n \in \N^*$, la matrice de centrage est :
    \begin{equation}
        \Jmat \;=\; \mat{I}_n - \frac{1}{n}\,\ones\ones^\top \;\in\; \R^{n\times n}
    \end{equation}
    où $\mat{I}_n$ est la matrice identité $n\times n$.
\end{definition}

\begin{proposition}[Propriétés de $\Jmat$]\label{prop:J}
    La matrice $\Jmat$ vérifie :
    \begin{enumerate}
        \item[\rm(i)] $\Jmat\ones = \vect{0}$
        \item[\rm(ii)] $\ones^\top\Jmat = \vect{0}^\top$
        \item[\rm(iii)] $\Jmat$ est symétrique : $\Jmat^\top = \Jmat$
        \item[\rm(iv)] $\Jmat$ est idempotente : $\Jmat^2 = \Jmat$
    \end{enumerate}
\end{proposition}

\begin{proof}
    \textbf{(i)} Par définition :
    \begin{align}
        \Jmat\ones
        &= \left(\mat{I}_n - \frac{1}{n}\,\ones\ones^\top\right)\ones \notag\\
        &= \mat{I}_n\ones - \frac{1}{n}\,\ones\,(\ones^\top\ones) \notag\\
        &= \ones - \frac{1}{n}\,\ones\cdot n
        \qquad\left[\text{car } \ones^\top\ones
            = \sum_{k=1}^n 1 = n\right]\notag\\
        &= \ones - \ones \;=\; \vect{0}
    \end{align}

    \textbf{(ii)} Découle de (i) par transposition : $\ones^\top\Jmat =
    (\Jmat^\top\ones)^\top = (\Jmat\ones)^\top = \vect{0}^\top$, en utilisant
    (iii).

    \textbf{(iii)} $\Jmat^\top = \mat{I}_n^\top - \frac{1}{n}(\ones\ones^\top)^\top
    = \mat{I}_n - \frac{1}{n}\ones\ones^\top = \Jmat$.

    \textbf{(iv)}
    $\Jmat^2 = \left(\mat{I} - \frac{1}{n}\ones\ones^\top\right)^2
    = \mat{I} - \frac{2}{n}\ones\ones^\top
    + \frac{1}{n^2}\ones(\ones^\top\ones)\ones^\top
    = \mat{I} - \frac{2}{n}\ones\ones^\top + \frac{1}{n}\ones\ones^\top
    = \mat{I} - \frac{1}{n}\ones\ones^\top = \Jmat$.
\end{proof}

\begin{remark}
    Appliqué à un vecteur $\vect{v} \in \R^n$, $\Jmat$ soustrait la moyenne :
    $(\Jmat\vect{v})_i = v_i - \bar{v}$ où $\bar{v} = \frac{1}{n}\sum_k v_k$.
    Appliqué à une matrice $\mat{M}$ sous la forme $\Jmat\mat{M}\Jmat$, il
    soustrait simultanément les moyennes de lignes et de colonnes et rajoute
    la moyenne globale : c'est le \textbf{double centrage}.
\end{remark}

\subsection{Théorème du double centrage}

\begin{theorem}[Double centrage]\label{thm:double_centering}
    Soit $\mat{X} \in \R^{n\times d}$ une configuration de points vérifiant
    $\sum_{i=1}^n \vect{x}_i = \vect{0}$ (centrage de la configuration).
    Soit $\Dsq \in \R^{n\times n}$ la matrice des distances euclidiennes au
    carré associée. Alors :
    \begin{equation}
        \boxed{-\frac{1}{2}\,\Jmat\Dsq\Jmat \;=\; \Gram \;=\; \mat{X}\mat{X}^\top}
    \end{equation}
\end{theorem}

\begin{proof}
    On part de la décomposition~\eqref{eq:gram_dist_matrix} :
    \[
        \Dsq = \vect{r}\ones^\top + \ones\vect{r}^\top - 2\Gram
    \]

    On calcule $-\frac{1}{2}\Jmat\Dsq\Jmat$ en distribuant :

    \begin{equation}
        -\frac{1}{2}\Jmat\Dsq\Jmat
        = -\frac{1}{2}\Jmat\!\left(\vect{r}\ones^\top\right)\!\Jmat
        -\frac{1}{2}\Jmat\!\left(\ones\vect{r}^\top\right)\!\Jmat
        + \Jmat\Gram\Jmat
        \label{eq:expand}
    \end{equation}

    \textbf{Annulation du premier terme.}
    En utilisant la propriété (ii) de la Proposition~\ref{prop:J},
    $\ones^\top\Jmat = \vect{0}^\top$, donc :
    \[
        \Jmat\!\left(\vect{r}\ones^\top\right)\!\Jmat
        = \Jmat\vect{r}\,\underbrace{(\ones^\top\Jmat)}_{=\,\vect{0}^\top}
        = \mat{0}
    \]

    \textbf{Annulation du second terme.}
    En utilisant la propriété (i), $\Jmat\ones = \vect{0}$, donc :
    \[
        \Jmat\!\left(\ones\vect{r}^\top\right)\!\Jmat
        = \underbrace{(\Jmat\ones)}_{=\,\vect{0}}\vect{r}^\top\Jmat
        = \mat{0}
    \]

    \textbf{Simplification du troisième terme.}
    On doit montrer que $\Jmat\Gram\Jmat = \Gram$. L'hypothèse de centrage
    $\sum_i \vect{x}_i = \vect{0}$ implique :
    \[
        (\mat{X}^\top\ones)_k = \sum_{i=1}^n x_{ik} = 0
        \quad\Longrightarrow\quad
        \mat{X}^\top\ones = \vect{0}
    \]

    Alors :
    \[
        \Gram\ones = \mat{X}\mat{X}^\top\ones = \mat{X}\,(\mat{X}^\top\ones)
        = \mat{X}\cdot\vect{0} = \vect{0}
    \]
    et par transposition $\ones^\top\Gram = \vect{0}^\top$. Ainsi :
    \begin{align}
        \Jmat\Gram\Jmat
        &= \left(\mat{I} - \frac{1}{n}\ones\ones^\top\right)\Gram
           \left(\mat{I} - \frac{1}{n}\ones\ones^\top\right) \notag\\
        &= \Gram
           - \frac{1}{n}\ones\underbrace{(\ones^\top\Gram)}_{=\,\vect{0}^\top}
           - \frac{1}{n}\underbrace{(\Gram\ones)}_{=\,\vect{0}}\ones^\top
           + \frac{1}{n^2}\ones\underbrace{(\ones^\top\Gram\ones)}_{=\,0}\ones^\top
        \notag\\
        &= \Gram
    \end{align}

    \textbf{Conclusion.} En reportant dans~\eqref{eq:expand} :
    \[
        -\frac{1}{2}\Jmat\Dsq\Jmat = 0 + 0 + \Gram = \Gram \qedhere
    \]
\end{proof}

\begin{remark}
    L'hypothèse $\sum_i \vect{x}_i = \vect{0}$ n'est pas restrictive : on
    peut toujours centrer une configuration en lui soustrayant son centroïde.
    Le double centrage de $\Dsq$ est exactement l'opération qui impose ce
    centrage implicitement. C'est pourquoi l'algorithme MDS n'a pas besoin
    d'imposer ce centrage a priori : il est automatiquement réalisé.
\end{remark}

\subsection{Récupération des coordonnées par décomposition propre}

\subsubsection{Propriétés de la matrice de Gram}

\begin{proposition}\label{prop:gram_psd}
    La matrice de Gram $\Gram = \mat{X}\mat{X}^\top$ est symétrique
    semi-définie positive (SSDP).
\end{proposition}

\begin{proof}
    \textbf{Symétrie :} $\Gram^\top = (\mat{X}\mat{X}^\top)^\top =
    \mat{X}\mat{X}^\top = \Gram$.

    \textbf{Semi-définie positive :} Pour tout $\vect{v} \in \R^n$ :
    \[
        \vect{v}^\top\Gram\vect{v}
        = \vect{v}^\top\mat{X}\mat{X}^\top\vect{v}
        = (\mat{X}^\top\vect{v})^\top(\mat{X}^\top\vect{v})
        = \norm{\mat{X}^\top\vect{v}}^2 \;\geq\; 0
    \]
\end{proof}

\subsubsection{Décomposition spectrale}

Le théorème spectral garantit que toute matrice symétrique réelle se
diagonalise dans une base orthonormée :

\begin{theorem}[Décomposition spectrale]\label{thm:spectral}
    Soit $\Gram \in \R^{n\times n}$ symétrique SSDP. Il existe une matrice
    orthogonale $\mat{V} \in \R^{n\times n}$ (i.e.\ $\mat{V}^\top\mat{V} =
    \mat{I}$) et une matrice diagonale $\mat{\Lambda} =
    \operatorname{diag}(\lambda_1, \ldots, \lambda_n)$ avec $\lambda_1 \geq
    \lambda_2 \geq \cdots \geq \lambda_n \geq 0$ telles que :
    \begin{equation}
        \Gram \;=\; \mat{V}\mat{\Lambda}\mat{V}^\top
    \end{equation}
    Les colonnes $\vect{v}_k$ de $\mat{V}$ sont les \textbf{vecteurs propres}
    de $\Gram$ et les $\lambda_k$ sont les \textbf{valeurs propres} associées,
    vérifiant $\Gram\vect{v}_k = \lambda_k\vect{v}_k$.
\end{theorem}

\subsubsection{Factorisation en racine de Gram}

On cherche $\mat{X} \in \R^{n\times 3}$ tel que $\mat{X}\mat{X}^\top = \Gram$.
En utilisant la décomposition spectrale :

\begin{align}
    \Gram &= \mat{V}\mat{\Lambda}\mat{V}^\top \notag\\
    &= \mat{V}\mat{\Lambda}^{1/2}\mat{\Lambda}^{1/2}\mat{V}^\top \notag\\
    &= \left(\mat{V}\mat{\Lambda}^{1/2}\right)
       \left(\mat{V}\mat{\Lambda}^{1/2}\right)^\top
\end{align}

où $\mat{\Lambda}^{1/2} = \operatorname{diag}(\sqrt{\lambda_1}, \ldots,
\sqrt{\lambda_n})$ est bien définie puisque $\lambda_k \geq 0$.

\begin{definition}[Solution MDS]
    En ne retenant que les 3 premières valeurs propres (correspondant à la
    dimension physique de l'espace), la solution MDS est :
    \begin{equation}
        \boxed{\mat{X} = \mat{V}_3\,\mat{\Lambda}_3^{1/2}}
    \end{equation}
    où $\mat{V}_3 \in \R^{n\times 3}$ contient les 3 premiers vecteurs propres
    et $\mat{\Lambda}_3 = \operatorname{diag}(\lambda_1, \lambda_2, \lambda_3)$.
\end{definition}

\begin{proposition}[Interprétation des valeurs propres]\label{prop:eigenvalues}
    Pour une configuration de points exactement en dimension 3, la matrice de
    Gram a exactement 3 valeurs propres non nulles : $\lambda_1 \geq \lambda_2
    \geq \lambda_3 > 0$ et $\lambda_4 = \cdots = \lambda_{17} = 0$.
    En présence de bruit de mesure, on observe $\lambda_4, \ldots,
    \lambda_{17} \approx 0 \ll \lambda_3$.
\end{proposition}

\begin{remark}[Signification physique des valeurs propres]
    La formule de reconstruction $\mat{X} = \mat{V}_3\mat{\Lambda}_3^{1/2}$
    implique que la $k$-ième colonne de $\mat{X}$ vaut $\sqrt{\lambda_k}\,\vect{v}_k$.
    Or cette colonne contient les coordonnées de tous les $n$ microphones sur
    le $k$-ième axe principal. On a donc directement :
    \[
        \lambda_k
        = \bigl\lVert \operatorname{col}_k(\mat{X}) \bigr\rVert^2
        = \bigl\lVert \sqrt{\lambda_k}\,\vect{v}_k \bigr\rVert^2
        = \lambda_k \underbrace{\norm{\vect{v}_k}^2}_{=\,1}
        = \sum_{i=1}^{n} x_{ik}^2
    \]
    $\lambda_k$ est donc la \textbf{somme des carrés des coordonnées} de tous
    les microphones sur l'axe $k$ — et non une distance directe. Puisque la
    configuration est centrée ($\sum_i x_{ik} = 0$), $\lambda_k/n$ est la
    \emph{variance} sur cet axe, et la distance typique d'un microphone au
    centre vaut :
    \[
        \sigma_k = \sqrt{\lambda_k / n} \quad [\text{mètres}]
    \]
    En acoustique de salle, on s'attend typiquement à
    $\lambda_1 \gg \lambda_2 > \lambda_3 \gg 0$,
    correspondant respectivement à la profondeur, la largeur et la hauteur de
    la scène : les microphones sont davantage étalés en profondeur qu'en
    hauteur.
\end{remark}

\subsubsection{Ambiguïté de rotation}

\begin{proposition}[Non-unicité de la solution]\label{prop:rotation_ambiguity}
    Si $\mat{X}$ est solution du problème MDS, alors pour toute matrice
    orthogonale $\mat{R} \in \R^{3\times 3}$ (rotation ou réflexion), la
    configuration $\mat{X}' = \mat{X}\mat{R}$ est également solution.
\end{proposition}

\begin{proof}
    Il suffit de vérifier que $\mat{X}'$ reproduit la même matrice de Gram :
    \[
        \mat{X}'\mat{X}'^\top
        = (\mat{X}\mat{R})(\mat{X}\mat{R})^\top
        = \mat{X}\mat{R}\mat{R}^\top\mat{X}^\top
        = \mat{X}\,\mat{I}_3\,\mat{X}^\top
        = \mat{X}\mat{X}^\top = \Gram
    \]
    La propriété $\mat{R}\mat{R}^\top = \mat{I}_3$ (orthogonalité) assure
    l'égalité.
\end{proof}

En conséquence, MDS fournit des positions correctes à une \emph{isométrie}
(rotation, réflexion, translation) près. L'algorithme de Kabsch
(Section~\ref{sec:kabsch}) est nécessaire pour aligner la solution dans le
repère physique de la salle de concert.


% ══════════════════════════════════════════════════════════════════════════════
\section{Phase 2 : Trilatération des microphones d'ambiance}
% ══════════════════════════════════════════════════════════════════════════════

\subsection{Formulation du problème}

Après la Phase 1, les positions $\vect{x}_1, \ldots, \vect{x}_{17} \in \R^3$
des microphones dédiés sont connues. Pour chaque microphone d'ambiance $k \in
\mathcal{A}$, on dispose de ses distances à tous les microphones dédiés :

\[
    d_{ki} = \norm{\vect{x}_k - \vect{x}_i}, \qquad i = 1,\ldots,17
\]

On cherche $\vect{x}_k = (x, y, z)^\top \in \R^3$. Ce problème est appelé
\textbf{trilatération} (analogue 3D du GPS).

\subsection{Système d'équations non-linéaires}

Le développement de $\norm{\vect{x}_k - \vect{x}_i}^2 = d_{ki}^2$ donne :

\begin{equation}
    \norm{\vect{x}_k}^2 - 2\,\vect{x}_i^\top\vect{x}_k + \norm{\vect{x}_i}^2
    = d_{ki}^2
    \qquad i = 1,\ldots,17
    \label{eq:trilateration_nonlinear}
\end{equation}

Ce système est non-linéaire en $\vect{x}_k$ à cause du terme
$\norm{\vect{x}_k}^2$.

\subsection{Linéarisation par soustraction}

Pour éliminer le terme non-linéaire $\norm{\vect{x}_k}^2$, on soustrait
l'équation $i=1$ de chacune des équations $i = 2, \ldots, 17$ :

\textbf{Équation pour $i = 1$ :}
\begin{equation}
    \norm{\vect{x}_k}^2 - 2\,\vect{x}_1^\top\vect{x}_k + \norm{\vect{x}_1}^2
    = d_{k1}^2
    \label{eq:eq1}
\end{equation}

\textbf{Équation pour $i \geq 2$ :}
\begin{equation}
    \norm{\vect{x}_k}^2 - 2\,\vect{x}_i^\top\vect{x}_k + \norm{\vect{x}_i}^2
    = d_{ki}^2
    \label{eq:eqi}
\end{equation}

En soustrayant~\eqref{eq:eq1} de~\eqref{eq:eqi}, le terme $\norm{\vect{x}_k}^2$
s'annule :

\begin{equation}
    -2\,(\vect{x}_i - \vect{x}_1)^\top\vect{x}_k
    + \norm{\vect{x}_i}^2 - \norm{\vect{x}_1}^2
    = d_{ki}^2 - d_{k1}^2
\end{equation}

En réarrangeant :

\begin{equation}
    -2\,(\vect{x}_i - \vect{x}_1)^\top\vect{x}_k
    = d_{ki}^2 - d_{k1}^2 - \norm{\vect{x}_i}^2 + \norm{\vect{x}_1}^2
    \label{eq:linear_eq}
\end{equation}

\subsection{Formulation matricielle par moindres carrés}

On définit pour $i = 2, \ldots, 17$ :

\begin{equation}
    a_i = -2\,(\vect{x}_i - \vect{x}_1)^\top \in \R^{1\times 3},
    \qquad
    b_i = d_{ki}^2 - d_{k1}^2 - \norm{\vect{x}_i}^2 + \norm{\vect{x}_1}^2 \in \R
\end{equation}

On obtient le système linéaire surdéterminé (16 équations, 3 inconnues) :

\begin{equation}
    \mat{A}\,\vect{x}_k = \vect{b}
    \label{eq:linear_system}
\end{equation}

avec :

\begin{equation}
    \mat{A} = \begin{pmatrix} a_2 \\ \vdots \\ a_{17} \end{pmatrix}
    = -2\begin{pmatrix}
        (\vect{x}_2 - \vect{x}_1)^\top \\
        \vdots \\
        (\vect{x}_{17} - \vect{x}_1)^\top
    \end{pmatrix}
    \in \R^{16\times 3},
    \qquad
    \vect{b} = \begin{pmatrix} b_2 \\ \vdots \\ b_{17} \end{pmatrix}
    \in \R^{16}
\end{equation}

Le système étant surdéterminé ($16 > 3$), on cherche la solution au sens des
\textbf{moindres carrés} :

\begin{equation}
    \boxed{\vect{x}_k^* = \underset{\vect{x} \in \R^3}{\arg\min}\;
    \norm{\mat{A}\vect{x} - \vect{b}}^2}
    \label{eq:lstsq}
\end{equation}

\begin{proposition}[Solution des moindres carrés]
    Si $\mat{A}$ est de rang plein (rang 3), la solution unique est :
    \begin{equation}
        \vect{x}_k^* = (\mat{A}^\top\mat{A})^{-1}\mat{A}^\top\vect{b}
    \end{equation}
    En pratique, on utilise la décomposition QR ou SVD de $\mat{A}$ pour
    une meilleure stabilité numérique (\texttt{numpy.linalg.lstsq}).
\end{proposition}

\begin{remark}[Pourquoi 16 équations améliorent la robustesse]
    Trois équations suffiraient théoriquement pour trois inconnues. En
    utilisant les 16 équations disponibles, la solution aux moindres carrés
    est un estimateur de \textbf{variance minimale} : les erreurs de mesure
    (bruit sur les TOA, réflexions parasites) sont moyennées, ce qui réduit
    l'erreur quadratique moyenne sur $\vect{x}_k^*$.
\end{remark}

\begin{remark}[Choix du point de référence]
    L'équation~\eqref{eq:linear_eq} utilise $\vect{x}_1$ comme point de
    référence. Ce choix est arbitraire ; en pratique on peut améliorer la
    stabilité numérique en choisissant le microphone dédié le plus proche du
    microphone d'ambiance à localiser, de sorte que $d_{k1}$ soit le plus
    petit possible.
\end{remark}

% ══════════════════════════════════════════════════════════════════════════════
\section{Orientation sémantique des axes sans Kabsch}
\label{sec:orientation}
% ══════════════════════════════════════════════════════════════════════════════

\subsection{Motivation}

La Proposition~\ref{prop:rotation_ambiguity} montre que la solution MDS est
définie à une isométrie près.  L'algorithme de Kabsch (Section~\ref{sec:kabsch})
résout ce problème par optimisation sur des points d'ancrage GPS, mais
requiert la connaissance préalable des coordonnées réelles de plusieurs
microphones.

Lorsque ces données ne sont pas disponibles, la \textbf{disposition connue
d'un orchestre symphonique} fournit des contraintes d'orientation suffisantes
pour lever les trois degrés de liberté (signe de chaque axe) et identifier
quelle colonne de $\mat{X}$ correspond à chaque direction physique.

\subsection{Contraintes orchestrales}

La disposition standard d'un orchestre symphonique impose les relations
suivantes entre les sections :

\begin{enumerate}
    \item \textbf{Profondeur} : les cordes (Vln, Vla, Vcl, CBasse) occupent
        les premiers rangs, les cuivres et percussions les rangs arrière.
        Le centroïde des cordes est donc \emph{devant} celui des cuivres.

    \item \textbf{Latéral} : les premiers violons sont à gauche du chef
        d'orchestre, les violoncelles à droite.  En convention SOFA
        ($90° =$ gauche), $\vect{x}_{\text{Vln}_1}$ a une composante latérale
        positive par rapport à $\vect{x}_{\text{Vcl}}$.

    \item \textbf{Hauteur} : tous les instruments sont sur la même scène,
        avec une faible variation verticale ($\leq 2\,\text{m}$).  L'axe de
        hauteur correspond donc au troisième axe principal (variance minimale).
\end{enumerate}

\subsection{Algorithme d'orientation}

Soit $\mat{X} \in \R^{n \times 3}$ la solution brute du cMDS, avec les
colonnes ordonnées par variance décroissante ($\lambda_1 \geq \lambda_2 \geq
\lambda_3$).  On définit les ensembles d'indices :

\begin{align*}
    \mathcal{S} &= \{\text{Vln}_1, \text{Vln}_2, \text{Vla}, \text{Vcl},
                     \text{CBasse}_1, \text{CBasse}_2\} \\
    \mathcal{B} &= \{\text{Cors}, \text{Trompettes}, \text{Trombones},
                     \text{Tuba}, \text{Timbales}, \text{GC}\}
\end{align*}

et les centroïdes de groupe :

\begin{equation}
    \bar{\vect{x}}_{\mathcal{S}} = \frac{1}{|\mathcal{S}|}
    \sum_{i \in \mathcal{S}} \vect{x}_i,
    \qquad
    \bar{\vect{x}}_{\mathcal{B}} = \frac{1}{|\mathcal{B}|}
    \sum_{i \in \mathcal{B}} \vect{x}_i
\end{equation}

\paragraph{Étape 1 — Identification des axes.}

\begin{align}
    k_{\rm prof}   &= \underset{k \in \{1,2,3\}}{\arg\max}\;
                      \bigl|(\bar{\vect{x}}_{\mathcal{B}}
                           - \bar{\vect{x}}_{\mathcal{S}})_k\bigr|
    \label{eq:depth_axis}\\[4pt]
    k_{\rm lat}    &= \underset{k \neq k_{\rm prof}}{\arg\max}\;
                      |(\vect{x}_{\text{Vln}_1} - \vect{x}_{\text{Vcl}})_k|
    \label{eq:lat_axis}\\[4pt]
    k_{\rm haut}   &= \{1,2,3\} \setminus \{k_{\rm prof},\, k_{\rm lat}\}
    \label{eq:height_axis}
\end{align}

L'axe de profondeur~\eqref{eq:depth_axis} est celui qui sépare le plus les
deux groupes sémantiques.  L'axe latéral~\eqref{eq:lat_axis} est identifié
parmi les deux restants par la séparation Vln$_1$/Vcl.

\paragraph{Étape 2 — Permutation.}

On réordonne les colonnes de $\mat{X}$ :
\begin{equation}
    \mat{X}' = \mat{X}[:,\; (k_{\rm prof},\; k_{\rm lat},\; k_{\rm haut})]
\end{equation}

\paragraph{Étape 3 — Correction des signes.}

\begin{align}
    \text{Si } (\bar{x}'_{\mathcal{B}})_1 < (\bar{x}'_{\mathcal{S}})_1
    &\;\Longrightarrow\; X'_{:,1} \leftarrow -X'_{:,1}
    \quad\text{(cuivres en }X > 0\text{)} \label{eq:sign_depth}\\[4pt]
    \text{Si } x'_{\text{Vln}_1, 2} < x'_{\text{Vcl}, 2}
    &\;\Longrightarrow\; X'_{:,2} \leftarrow -X'_{:,2}
    \quad\text{(Vln}_1\text{ en }Y > 0\text{)} \label{eq:sign_lat}
\end{align}

\begin{remark}[Limites de l'approche]
    Cette méthode lève les ambiguïtés de signe ($\pm 1$ sur chaque axe) et
    identifie la correspondance axe MDS $\leftrightarrow$ direction physique,
    mais ne corrige pas l'angle de rotation dans le plan horizontal (angle de
    la scène par rapport à l'axe profondeur).  Pour un recalage métrique
    complet, l'algorithme de Kabsch (Section~\ref{sec:kabsch}) reste nécessaire
    dès lors que les coordonnées réelles d'au moins trois microphones sont
    connues.
\end{remark}

% ══════════════════════════════════════════════════════════════════════════════
\section{Alignement de coordonnées : algorithme de Kabsch}
\label{sec:kabsch}
% ══════════════════════════════════════════════════════════════════════════════

\subsection{Motivation}

La Proposition~\ref{prop:rotation_ambiguity} montre que MDS renvoie une
configuration définie à une isométrie près.  Lorsque les coordonnées réelles
de $p \geq 3$ points d'ancrage $\vect{q}_1, \ldots, \vect{q}_p \in \R^3$
sont connues (plan de scène, positions GPS), on cherche la rotation
$\mat{R} \in \R^{3\times3}$ et la translation $\vect{t} \in \R^3$ qui
minimisent l'erreur de recalage :

\begin{equation}
    (\mat{R}^*, \vect{t}^*) = \underset{\substack{\mat{R}^\top\mat{R}=\mat{I}\\\det(\mat{R})=1}}{\arg\min}
    \sum_{i=1}^p \norm{\mat{R}\,\vect{p}_i + \vect{t} - \vect{q}_i}^2
\end{equation}

où $\vect{p}_i$ sont les positions estimées pour ces mêmes points.

\subsection{Algorithme de Kabsch}

\begin{enumerate}
    \item \textbf{Centrage :}
    \[
        \bar{\vect{p}} = \frac{1}{p}\sum_i\vect{p}_i, \quad
        \bar{\vect{q}} = \frac{1}{p}\sum_i\vect{q}_i, \quad
        \tilde{\vect{p}}_i = \vect{p}_i - \bar{\vect{p}}, \quad
        \tilde{\vect{q}}_i = \vect{q}_i - \bar{\vect{q}}
    \]

    \item \textbf{Matrice de covariance croisée :}
    \[
        \mat{H} = \sum_{i=1}^p \tilde{\vect{p}}_i\,\tilde{\vect{q}}_i^\top
        = \mat{P}_c^\top\mat{Q}_c \;\in\; \R^{3\times 3}
    \]

    \item \textbf{SVD :} $\mat{H} = \mat{U}\mat{S}\mat{V}^\top$

    \item \textbf{Rotation optimale :}
    \[
        \mat{R}^* = \mat{V}\,\mat{U}^\top
    \]
    Si $\det(\mat{R}^*) = -1$ (réflexion), on corrige :
    $\mat{V}_{:,3} \leftarrow -\mat{V}_{:,3}$, puis $\mat{R}^* = \mat{V}\mat{U}^\top$.

    \item \textbf{Translation :} $\vect{t}^* = \bar{\vect{q}} - \mat{R}^*\bar{\vect{p}}$

    \item \textbf{Application à toutes les positions :}
    \[
        \vect{x}_j^{\text{aligné}} = \mat{R}^*\,\vect{x}_j + \vect{t}^*
    \]
\end{enumerate}

% ══════════════════════════════════════════════════════════════════════════════
\section{Conversion en coordonnées sphériques}
% ══════════════════════════════════════════════════════════════════════════════

Une fois les positions cartésiennes $\vect{x}_j = (x_j, y_j, z_j)^\top$
disponibles dans le repère de la salle (axe $x$ vers la scène, axe $y$ vers
la gauche, axe $z$ vers le haut), on les convertit en coordonnées sphériques
pour l'interpolateur HRTF :

\begin{align}
    r_j &= \norm{\vect{x}_j} \\
    \theta_j &= \operatorname{atan2}(y_j,\, x_j)
        \cdot \frac{180}{\pi} \quad [\text{azimut, degrés}]\\
    \phi_j &= \arcsin\!\left(\frac{z_j}{r_j}\right)
        \cdot \frac{180}{\pi} \quad [\text{élévation, degrés}]
\end{align}

La convention SOFA CCW est appliquée : $0°$ = devant, $90°$ = gauche,
$180°$ = derrière, $270°$ = droite.

% ══════════════════════════════════════════════════════════════════════════════
\section*{Conclusion}
% ══════════════════════════════════════════════════════════════════════════════

Ce document présente la chaîne mathématique complète de localisation des 23
microphones à partir des données TOA.

Deux corrections robustifient l'extraction des distances.
D'abord, le critère de sélection du microphone dédié utilise
$\arg\min(s_{ij}/a_{ij})$ plutôt que $\arg\min(s_{ij})$ seul, afin d'éliminer
les artefacts de déclenchement à faible amplitude.  Ensuite, la distance entre
deux micros dédiés est estimée par $\min(s_{ij}, s_{ji})/f_s \cdot c$ plutôt
que par $s_{ij}/f_s \cdot c$ : prendre le minimum des deux sens de mesure
fournit une borne supérieure plus serrée sur le trajet direct, en évitant les
réflexions tardives qui peuvent dominer le maximum de la RIR.

La Phase 1 exploite la structure bilinéaire des distances euclidiennes au carré
pour récupérer la matrice de Gram par double centrage, puis la décompose par
diagonalisation pour obtenir les coordonnées 3D des 17 microphones dédiés.
L'orientation des axes est résolue sans points d'ancrage GPS, par contraintes
sémantiques sur la disposition orchestrale (profondeur cordes/cuivres, latéral
Vln$_1$/Vcl).

La Phase 2 linéarise le problème de trilatération par soustraction d'équations,
permettant une résolution robuste par moindres carrés pour les 5 microphones
d'ambiance.  L'algorithme de Kabsch offre un recalage métrique complet lorsque
les coordonnées réelles de points d'ancrage sont disponibles.

\end{document}